\section{Methodology}
\label{sec:methodology}

\subsection{System Overview}

GrabDrop employs a two-stage gesture detection pipeline shared across
both the Neural Network (TCN) and Legacy (rule-based) backends
(Figure~\ref{fig:two-stage}):

\begin{enumerate}[leftmargin=*]
  \item \textbf{IDLE stage} ($\sim$10\,FPS): Low-power scanning for
    hand presence. A sliding window of 10~frames tracks hand detection
    counts; when $\geq 8$ frames contain a detected hand, the system
    transitions to the WAKEUP stage.
  \item \textbf{WAKEUP stage} ($\sim$30\,FPS): High-precision gesture
    classification. Runs for at most 2~seconds before timing out and
    returning to IDLE.
\end{enumerate}

\begin{figure}[ht]
  \centering
  \includegraphics[width=0.85\linewidth]{two-stage-pipeline}
  \caption{Two-stage detection pipeline: IDLE (low power) $\to$ WAKEUP
    (high precision). The system spends most time in IDLE, consuming
    minimal CPU.}
  \label{fig:two-stage}
\end{figure}

\subsection{Neural Network Backend: Temporal Convolutional Network}

The core gesture classifier is a Temporal Convolutional Network
(TCN)~\cite{bai2018empirical, lea2017temporal} that classifies a
30-frame sliding window of 144-dimensional feature vectors into one
of five classes.

\subsubsection{Causal Dilated Convolutions}

The TCN uses \emph{causal} 1D convolutions: the output at time $t$
depends only on inputs at times $\leq t$, which is essential for
real-time inference where future frames are unavailable. \emph{Dilated}
convolutions skip intermediate time steps to expand the receptive field
without increasing parameters:

\begin{equation}
  y(t) = \sum_{i=0}^{k-1} w(i) \cdot x(t - d \cdot i),
  \label{eq:dilated_conv}
\end{equation}
where $k$ is the kernel size, $d$ is the dilation factor, and $w$
denotes learned filter weights. With dilations $\{1, 2, 4\}$, the
network captures local, medium-range, and long-range temporal patterns
respectively.

\subsubsection{Residual Blocks}

Each processing stage uses residual blocks with the structure:
\[
  \text{output} = \text{ReLU}\!\bigl(\mathcal{F}(x) + x\bigr),
\]
where $\mathcal{F}$ denotes two stacked causal convolutions with batch
normalization and dropout ($p=0.15$). The skip connection $+\,x$
facilitates gradient flow and enables the block to learn incremental
refinements. When the channel dimension changes, a $1\!\times\!1$
convolution serves as the skip projection.

\subsubsection{Model Architecture}

The full GestureTCN architecture comprises:

\begin{enumerate}[leftmargin=*]
  \item \textbf{Stem}: $1\!\times\!1$ convolution projecting 144 input
    features to 48 channels, followed by batch normalization and ReLU.
  \item \textbf{Block~1}: Residual block, $k\!=\!3$, dilation $d\!=\!1$
    (48 channels, local patterns).
  \item \textbf{Block~2}: Residual block, $k\!=\!3$, $d\!=\!2$
    (48 channels, medium-range patterns).
  \item \textbf{Block~3}: Channel block, $k\!=\!3$, $d\!=\!4$
    (48$\to$64 channels, long-range patterns).
  \item \textbf{Block~4}: Residual block, $k\!=\!3$, $d\!=\!1$
    (64 channels, refinement).
  \item \textbf{Pooling}: Adaptive average pooling across the time
    dimension, producing a single 64-dim vector.
  \item \textbf{Classification head}: Two fully connected layers
    ($64 \to 32 \to 5$) with ReLU and dropout.
\end{enumerate}

The model contains only \textbf{87,077 parameters} (0.34\,MB),
making it suitable for mobile and embedded deployment. Weights are
initialized using Kaiming normal initialization~\cite{he2015delving}.

\subsection{Legacy Backend: Rule-Based Detection}

As a fallback, GrabDrop includes a rule-based detection mode that
requires no trained model. For each frame, it computes a finger curl
ratio:
\begin{equation}
  r_i = \frac{\|\text{tip}_i - \text{wrist}\|}
             {\|\text{mcp}_i - \text{wrist}\|},
  \label{eq:curl_ratio}
\end{equation}
where $\text{tip}_i$ and $\text{mcp}_i$ are the fingertip and
metacarpophalangeal joint of finger $i$ (excluding the thumb). A finger
is classified as \emph{extended} if $r_i > 1.3$ and \emph{curled} if
$r_i < 0.9$. A hand is classified as \textsc{palm} when $\geq 3$
fingers are extended, and \textsc{fist} when $\geq 3$ are curled.

In the WAKEUP stage, 8~consecutive frames of the target state
(e.g., \textsc{fist} following a \textsc{palm} trigger) confirm a
\texttt{grab} gesture. Swipe detection uses frame-to-frame wrist
$y$-velocity and cumulative displacement against configurable
thresholds.

If the Neural Network backend is selected but fails to load (e.g.,
missing ONNX file), the system automatically falls back to Legacy
mode.
